\documentclass[11pt,a4paper]{article}
% The point model in Fortran (6 Oct 2026): a note for the partner group and Keio.
\usepackage[margin=22mm,top=20mm,bottom=20mm]{geometry}
\usepackage[T1]{fontenc}
\usepackage{mathptmx}
\usepackage{amsmath}
\usepackage{booktabs}
\usepackage{enumitem}
\usepackage{xcolor}
\usepackage[hidelinks]{hyperref}
\definecolor{c3}{RGB}{30,80,160}
\setlength{\parindent}{0pt}
\setlength{\parskip}{4pt}
\setlist{nosep,leftmargin=1.2em}
\newcommand{\sect}[1]{\par\vspace{6pt}{\color{c3}\large\bfseries #1}\par\vspace{2pt}}
\newcommand{\file}[1]{\texttt{\small #1}}

\begin{document}
\sloppy
{\Large\bfseries The point model in Fortran: coupling without the material server}\par
\vspace{3pt}
Keisuke Nishioka, IKM, Leibniz Universit\"at Hannover, 6 October 2026\par
\vspace{2pt}\hrule

\sect{Why}
In the coupled runs the field of Klempt et al.\ (2024) sets the amount of
biofilm $\phi$ and the point model of Klempt et al.\ (2026) sets its
composition at every integration point. Until now the point model ran in
Python: the material routine sent the state of each integration point to a
material server through a socket and waited for the answer. This call took
about 100 times as long as the rest of a run. Two species on 2880 elements
needed 1.5\,h in Abaqus, one species on $40^3$ elements without the server
16\,min.

I therefore wrote the point model in Fortran (\file{ecology\_native.f}). It
replaces the Python bridge in the build and keeps its interface, so the material
routines of ANSYS and Abaqus use it without any change. No server, no socket
and no C layer are needed.

\sect{What it computes}
The state of one point is
$\mathbf g=(\phi_1,\dots,\phi_5,\ \phi_0,\ \psi_1,\dots,\psi_5,\ \gamma)$, with
$n\le5$ active species (unused species stay exactly zero). One coupling step of
length $s\,\Delta t$ is divided into sub-steps of at most $10^{-4}$. In each
sub-step the twelve residuals $\mathbf R(\mathbf g_{k+1};\mathbf g_k)=\mathbf 0$
of the point model (Klempt et al.\ 2026, Eq.~16--18, divided by $\eta_i$) are
solved by six Newton iterations,
\[
  \mathbf g \leftarrow \mathbf g-\Big(\frac{\partial\mathbf R}{\partial\mathbf g}\Big)^{-1}\mathbf R ,
\]
with the state clipped to its admissible range before and after each iteration.
The Jacobian is the exact derivative, written out by hand; the linear system is
solved by LU decomposition with partial pivoting. This is the same scheme as the
Python reference, which computes the Jacobian by automatic differentiation.

The constants of a case ($n$, $c^*$, $\alpha^*$, $\eta_i$ and the interaction
matrix) are read once from a small text file, written by
\file{write\_eco\_cfg.py <case>} from the same case table as the Python server.
As in the server, the interaction matrix in the input deck is checked against
the case. The local nutrient enters as $c^*=c^*_0\,\min(\max(c/c_{\rm ref},0),1)$.

\sect{Checks}
\begin{center}
\begin{tabular}{p{0.52\linewidth}p{0.38\linewidth}}
\toprule
Check & Result \\
\midrule
Single points against the Python reference: 16 states, two-species cases 3 and 6,
four-species case 1 and the five-species model, several sub-step counts and
nutrient levels (\file{tests/test\_ecology\_native.py}) & largest relative difference $8\times10^{-15}$ \\
ANSYS, the partner's element, $8^3$ elements, two species, same deck as a run
with the Python server & stresses and composition identical in every printed
digit; trace within $1.6\times10^{-12}$ \\
Abaqus, $8^3$ elements, two species with the nutrient, same deck as a run with
the Python server & printed values identical; trace of 8213 point steps within
$6.7\times10^{-12}$ \\
\bottomrule
\end{tabular}
\end{center}
The remaining differences are round-off: the Fortran Jacobian and the
automatically differentiated one sum their terms in a different order.

\sect{Speed}
\begin{center}
\begin{tabular}{lrr}
\toprule
Run & Python server & Fortran \\
\midrule
ANSYS, partner element, $8^3$, two species & 156\,s & 36\,s \\
Abaqus, $8^3$, two species with the nutrient & 5.2\,min (4 CPUs) & 0.6\,min (1 CPU) \\
Abaqus, $16^3$, two species & 10\,min for 30\,\% of the run & 1\,min (whole run) \\
\bottomrule
\end{tabular}
\end{center}
The point model now costs about as much as the mechanics. Runs that were out of
reach with the server, such as two species on $32^3$ elements, take hours
instead of days.

\sect{How to use it}
\begin{itemize}
\item ANSYS: build the custom executable with \file{ecology\_native.f} in place of
  the Python bridge (\file{usermat\_py\_hook.f}); the material routine and the
  call-site code stay as they are. Before starting ANSYS, set the environment
  variable \file{BIOFILM\_ECO\_CASE} to the case file. One process per run.
\item Abaqus: \file{run\_comp.ps1 -Native} builds and links the same file into the
  user subroutine; the case file is read once at the start of the analysis.
\item The number of species in the material input must match the case file, as
  it must match the server's case.
\end{itemize}

\sect{Limits}
\begin{itemize}
\item Cases with a nutrient level or antibiotic level that changes in time
  (four-species cases 3 and 4 of Klempt et al.\ 2026) are not supported; the
  Python server does not support them either.
\item The Hill gate of this repository's five-species model is off, as in all
  current runs.
\item The number of Newton iterations is fixed at six, as in the reference. Where a
  sub-step is too long for six iterations to converge, both versions end at the
  clip bounds and can differ there; the sub-steps of the composition runs
  converge.
\end{itemize}

\vspace{4pt}
{\small\raggedright Source: \file{ansys\_usermat/coupling/ecology\_native.f}; checks and
timings: \file{abaqus\_composition/README.md}; repository
\href{https://github.com/keisuke58/pde-fem-biofilm}{keisuke58/pde-fem-biofilm}.}
\end{document}
